\documentclass{article}
\usepackage{iclr2027_conference,times}
\usepackage[T1]{fontenc}
\usepackage{amsmath,amssymb,booktabs,tabularx,microtype}
\usepackage{hyperref}
\usepackage{url}
\hypersetup{hidelinks}
\usepackage[section]{placeins}
\usepackage{tikz}
\usetikzlibrary{arrows.meta,positioning}
\setlength{\emergencystretch}{2em}
\iclrfinalcopy

\newcommand{\EV}{\mathrm{EV}}
\newcommand{\arm}[1]{\texttt{#1}}

\title{From CoSpRo to Cross-Fitted Concept Distillation:\\
A Research Log of the \texttt{concept-factors} Branch}
\author{Spur\_SpLiCE project\\
Branch \texttt{concept-factors}, 26 September to 2 October 2026}

\begin{document}
\maketitle

\begin{abstract}
This document records how the \texttt{concept-factors} branch left the CoSpRo method and arrived at
cross-fitted concept distillation. CoSpRo builds a teacher graph from concept projections of a frozen
CLIP/SpLiCE model and lifts Waterbirds worst-group accuracy (WGA) by four points over SimCLR, while on
MetaShift it stays below SimCLR. We replaced its prior about which concept is spurious with a
symmetric goal: keep strongly correlated concept factors on separate linear directions of the student,
so that a group-balanced linear probe can pick the side the task needs. Six rounds of experiments
followed. Distilling SpLiCE factors into a SimCLR ResNet doubles WGA on Spur-CIFAR10 from the first
attempt and leaves MetaShift unchanged through reweighting, whitening and per-factor contrastive
blocks. Shuffled-target controls showed that every per-image auxiliary loss is met by memorization on
MetaShift's 1,700 images; a ridge regression fitted in closed form on one half of each batch
(cross-fitting) removed memorization in the head and isolated the concept content as the source of the
Spur-CIFAR10 gain. A revised concept grouping, learnability bounds from CLIP features and a hold-out of
images from the concept loss then showed that the student memorizes in its backbone on MetaShift and
generalizes on Spur-CIFAR10. The configuration that came out of these diagnostics, meaning-grouped
factors distilled by the cross-fit at weight 10, reaches 55.3\% validation WGA on Waterbirds against
45.8\% for SimCLR and 42.8\% for its shuffled control, the best Waterbirds result of the project from
random initialization. It is now frozen and runs unchanged on CelebA, Spur-CIFAR10 and further
Waterbirds seeds.
\end{abstract}

\section{Overview}
\label{sec:overview}

Table~\ref{tab:timeline} lists the stages of the branch. Each stage answered a question the previous
one raised, and most answers were negative. The negative answers matter as much as the final positive
one, because each of them removed an explanation for the MetaShift failure and left the next one in
view: the choice of target factors, then the loss, then memorization in the head, then the grouping,
then the frequency of the factors and finally memorization in the backbone.

\begin{table}[h]
\centering
\caption{Stages of the branch. Dates are 2026; WGA is validation worst-group accuracy.}
\label{tab:timeline}
\small
\begin{tabularx}{\linewidth}{l l X}
\toprule
Date & Stage & Outcome \\
\midrule
before 26.09 & CoSpRo & +3.8 WGA on Waterbirds test; below SimCLR on MetaShift \\
26.09 & Concept factors, F1 and F2 & F2 doubles Spur-CIFAR10 WGA; MetaShift unchanged \\
27.09 & Rounds 1 and 2 on MetaShift & loss weight, whitening, atypicality weights: majority groups rise, minority groups stay near 50 \\
27--28.09 & Round 3 & balancing weights and concept blocks (F3); shuffled controls fit as well as real targets \\
28.09 & Cross-fitted F2 (xfit) & first version diverges; the repaired xfit matches F2 on Spur-CIFAR10 and its control falls to SimCLR \\
28.09 & Concept grouping revisited & synonyms rarely co-activate; meaning grouping by raw text embeddings \\
29.09 & Learnability bounds & rare factors are unpredictable from a half batch; factors from 5\% \\
29--30.09 & Loss weight 3 and 10 & training-image scores exceed CLIP's three times; WGA flat \\
30.09 & Unseen-image diagnostic & MetaShift memorizes, Spur-CIFAR10 generalizes \\
01--02.10 & Waterbirds & 55.3 WGA against 45.8 for SimCLR; configuration frozen \\
\bottomrule
\end{tabularx}
\end{table}

\section{Setting and protocol}
\label{sec:protocol}

All runs share the CoSpRo protocol. A ResNet-18 trains from random initialization with SimCLR
\citep{chen2020simclr} for 500 epochs (CelebA: 250) at batch size 128, temperature 0.05, learning rate
0.01 and float16 mixed precision on one V100. Every 25 epochs a logistic probe is fitted on a
group-balanced training subset and evaluated on the validation split; we report the mean of the last
four probes and the standard deviation over seeds unless a table says otherwise. Worst-group accuracy
is the lowest accuracy over the groups formed by class and spurious attribute \citep{sagawa2020groupdro}.
Concept factors come from SpLiCE \citep{bhalla2024splice} decompositions of OpenCLIP ViT-B/32
embeddings \citep{radford2021clip} over the LAION (10,000 words) or Open Images V7 (20,931 words)
vocabularies. Pretraining uses no class or group label.

\begin{table}[h]
\centering
\caption{Datasets. MetaShift and Spur-CIFAR10 are development datasets, on which settings are chosen;
Waterbirds and CelebA are held out for the final evaluation.}
\label{tab:datasets}
\small
\begin{tabularx}{\linewidth}{l r X}
\toprule
Dataset & Train images & Spurious structure \\
\midrule
MetaShift cats and dogs \citep{liang2022metashift} & 1,700 & cats indoors and dogs outdoors in 88\% of the images; 72 validation images per group \\
Spur-CIFAR10 & 45,000 & CIFAR-10 classes with a class-correlated coloured line; 100 groups \\
Waterbirds \citep{sagawa2020groupdro} & 4,795 & bird type correlated with land or water background \\
CelebA \citep{liu2015celeba} & 162,770 & hair colour correlated with gender \\
\bottomrule
\end{tabularx}
\end{table}

Two protocol rules shaped the decisions below. The development datasets carry all tuning; Waterbirds
and CelebA receive frozen configurations, so the label-free claim holds for them. MetaShift WGA is
noisy: SimCLR over ten seeds scores $44.6\pm6.9$, so two-seed differences below five points carry
little evidence and we leaned on mechanistic diagnostics more than on WGA there.

\section{Starting point: CoSpRo}
\label{sec:cospro}

CoSpRo removes the subspace of one concept group from CLIP embeddings, keeps image pairs that become
closer, filters them by activation contrast, residual concept agreement and synthetic null groups and
transfers the resulting sparse graph to a SimCLR student through graph-aware batches and a relational
KL loss. On Waterbirds it reaches $50.64\pm1.97$ test WGA against $46.82\pm3.65$ for SimCLR over four
seeds. Two observations motivated the branch.

\paragraph{The gain does not transfer.} On MetaShift, ten seeds give $42.9\pm2.7$ WGA for CoSpRo
against $44.6\pm6.9$ for SimCLR, and the LAION variant reaches $47.7\pm4.3$, within the SimCLR spread.
On Spur-CIFAR10 one seed gives 19.6 against 15.6.

\paragraph{The method carries a prior.} Projection relations help when removing a concept reveals a
same-class pair across contexts, which assumes that the spurious attribute is a removable context
concept. Label-free alternatives carry other priors: LateTVG \citep{hamidieh2024views} and
learning-speed methods \citep{zhu2023lassl} assume that the spurious feature is the easy one. The
CoSpRo graph also contains wrong-class edges, and background stays linearly decodable at 91\%.

\section{Concept factors: the idea}
\label{sec:idea}

Deep feature reweighting \citep{kirichenko2023dfr} shows that a linear head retrained on group-balanced
data recovers worst-group accuracy whenever the representation encodes the core feature apart from the
spurious one. Our probe is exactly such a head. The SSL failure to prevent is therefore the fusion of
two correlated features into one direction, or the suppression of the harder feature by the easier one,
which contrastive losses are known to produce \citep{chen2021intriguing,robinson2021shortcut}. Keeping
every strongly correlated pair of concepts apart serves both sides of the pair, and the balanced probe
chooses the side the task needs. No concept is declared spurious.

SpLiCE supplies the ingredients without labels. A \emph{factor} is a concept group of the CoSpRo
grouping stage; its value on image $i$ is the summed SpLiCE code $a_{ik}=\sum_{c\in G_k} c_{ic}$. Factors
outside a frequency band (2\% to 90\% of the images by default) are dropped. Two mechanisms used them.

\paragraph{F1, concept-conditioned batches.} A share of the batches holds only images that show one
factor of a strongly correlated pair (phi coefficient of presences). Inside such a batch the factor
tells no image apart, so InfoNCE has to separate the images by everything else, including the partner
factor. This is conditional contrastive learning \citep{tsai2021conditional} with a discovered factor
in place of an annotated attribute.

\paragraph{F2, factor distillation.} A linear head $W$ on the backbone feature $h_i$ predicts the
factor targets $t_i$ of both views,
\begin{equation}
\mathcal{L}=\mathcal{L}_{\mathrm{SimCLR}}+\lambda\,\frac{1}{2B}\sum_{i,v}\lVert W h_i^{(v)}-t_i\rVert^2/K .
\end{equation}
The targets are standardized factor values or their ZCA-whitened version \citep{kessy2018whitening},
which leaves each factor the part the other factors do not explain and so is large exactly on the
images that break a correlation. The head is linear so that the factors must be linearly readable, as
the probe reads them. A \emph{shuffled} control assigns each image the targets of another image: same
statistics, no image content.

\section{Rounds 1 to 3: distillation helps where data is plentiful}
\label{sec:rounds}

\subsection{Round 1}
The first inspection of the factors showed two defects, both fixed before training. Taking the 64 most
frequent groups dropped the scene concepts of MetaShift and most line colours of Spur-CIFAR10, and the
top correlated pairs were two names for one thing (cat breeds, lorry and truck). We kept every group in
the band and merged groups whose CLIP image responses correlate above 0.9.

\begin{table}[h]
\centering
\caption{Round 1, seeds 1 and 2. EV is the explained variance of the targets by the trained head.}
\label{tab:round1}
\small
\begin{tabular}{l l r r r}
\toprule
Dataset & Arm & WGA & Acc. & EV \\
\midrule
Spur-CIFAR10 & SimCLR & $15.6\pm1.0$ & 59.5 & \\
 & F2, standardized & $32.9\pm1.6$ & 69.5 & 0.62 \\
 & F2, shuffled & $24.2\pm0.4$ & 58.6 & 0.02 \\
\midrule
MetaShift & SimCLR & $47.7\pm4.2$ & 55.8 & \\
 & F2, standardized & $43.2\pm1.2$ & 58.2 & 0.12 \\
 & F2, shuffled & $42.2\pm4.7$ & 51.0 & 0.06 \\
\bottomrule
\end{tabular}
\end{table}

On Spur-CIFAR10 the gain was large and content-specific: real targets beat the shuffled control by
nine points of WGA and eleven of accuracy (Table~\ref{tab:round1}). The control itself added nine points
over SimCLR at equal accuracy, so the auxiliary regression also helps as such; Section~\ref{sec:xfit}
removes this ambiguity. On MetaShift the head explained 12\% of the factor variance and the arms
differed by less than the seed noise. We counted Spur-CIFAR10 as a success and MetaShift as an
unexplained failure.

\subsection{Round 2: pointing the loss at the minority images}
Every MetaShift arm with real targets raised average accuracy, while the minority groups (cats
outdoors and dogs indoors) stayed near 45 to 55\%. Our reading was that standardized targets reproduce
the dataset's correlations, and several Open Images concepts fuse class and context (``Cat bed'', ``Dog
walking''). Round 2 therefore tested ways to point F2 at the images that break the correlations without
labels: loss weights 3 and 10, whitened targets and per-image \emph{atypicality} weights (the mean
squared whitened factor value).

\begin{table}[h]
\centering
\caption{MetaShift, seeds 1 to 4 unless noted. Group accuracies in the order cat-indoor, cat-outdoor,
dog-indoor, dog-outdoor.}
\label{tab:round2}
\small
\begin{tabular}{l r r c r}
\toprule
Arm & WGA & Acc. & Groups & EV \\
\midrule
SimCLR & $44.8\pm9.3$ & 54.8 & 55 / 49 / 53 / 62 & \\
F2, weight 3 & $46.6\pm5.7$ & 61.5 & 70 / 47 / 54 / 74 & 0.41 \\
F2, weight 3, shuffled & $43.3\pm8.1$ & 54.8 & 57 / 45 / 53 / 65 & 0.33 \\
F2, weight 10 (seeds 1, 2) & $47.0\pm7.1$ & 59.7 & 62 / 46 / 53 / 75 & 0.85 \\
F2, whitened, weight 3 & $47.9\pm6.4$ & 56.9 & 57 / 51 / 55 / 65 & 0.31 \\
F2, weight 3, atypicality & $50.5\pm3.8$ & 58.8 & 59 / 52 / 55 / 70 & 0.61 \\
\bottomrule
\end{tabular}
\end{table}

Whitened targets and atypicality weights narrowed the gap between groups by lowering the majority gain;
none of them lifted a minority group above the SimCLR spread (Table~\ref{tab:round2}). We read this as
the regression learning through the co-occurrence itself: in 88\% of the images the indoor features
predict the cat factors, so one fused ``cat and indoor'' direction serves the loss.

\subsection{Round 3: making the co-occurrence useless}
Two mechanisms targeted the fused direction directly.
\emph{Balancing weights} minimize the mean squared correlation of the factor presences under per-image
weights with an entropy penalty, so a co-occurrence stops paying off in the weighted regression; this
is sample reweighting for independence as in stable learning \citep{zhang2021stable}.
\emph{F3 concept blocks} give every factor its own small linear block and apply a supervised contrastive
loss \citep{khosla2020supcon} inside the block, with images that share the factor as positives and a
pair weight that grows with how much the two images' other factors differ. A fused direction cannot
separate ``cat on sofa'' from ``dog on sofa'' inside the cat block, so the backbone needs a cat
direction that survives the context; pairs that share one factor and differ elsewhere identify that
factor in the weak-supervision sense of \citet{locatello2020weakly} and \citet{yao2024multiview}.

\begin{table}[h]
\centering
\caption{Round 3 on MetaShift, seeds 1 to 4.}
\label{tab:round3}
\small
\begin{tabular}{l r r c}
\toprule
Arm & WGA & Acc. & Groups \\
\midrule
F2, weight 3, balancing weights & $44.2\pm8.1$ & 59.0 & 64 / 44 / 55 / 72 \\
F3 concept blocks & $43.8\pm4.1$ & 60.0 & 69 / 45 / 51 / 75 \\
F3, all positive pairs alike & $49.7\pm2.5$ & 62.6 & 73 / 50 / 52 / 75 \\
F3, shuffled concept sets & $46.4\pm3.5$ & 56.2 & 56 / 48 / 52 / 68 \\
F3 and balanced F2 & $49.4\pm3.8$ & 64.4 & 74 / 52 / 55 / 78 \\
\bottomrule
\end{tabular}
\end{table}

No arm lifted a minority group (Table~\ref{tab:round3}). The controls explained why. F2 at weight 3
explained 0.41 of the real target variance and 0.33 of the shuffled one, and the block loss ended at
3.34 with real concepts and 3.37 with shuffled ones. On 1,700 images a network can fit per-image targets
through the features SimCLR learns anyway to tell images apart, and random targets are fitted almost
as well as real ones, as \citet{zhang2017rethinking} and \citet{arpit2017closer} showed for labels. Any
auxiliary loss attached to single images is then met without learning concepts. This was the first
mechanistic finding of the branch, and it redirected the work from the target design to the
estimator.

\section{Cross-fitting: no memorization in the head}
\label{sec:xfit}

The cross-fitted F2 (xfit) replaces the trained head by a ridge regression fitted in closed form on one
half $A$ of the batch and evaluated on the other half $B$, then the reverse. With unit-norm, centred
features $H$, targets $T$ and per-image weights $D$, the kernel form
\begin{equation}
\hat T_B = H_B H_A^\top \bigl(D_A H_A H_A^\top+\rho I\bigr)^{-1} D_A T_A
\end{equation}
needs one $n\times n$ solve and is differentiable, as in the meta-learning heads R2D2
\citep{bertinetto2019r2d2} and MetaOptNet \citep{lee2019metaoptnet}. The loss is the held-out error on
$B$. Both views of an image stay in one half. A feature that only identifies an image carries nothing
from $A$ to $B$, so only concept directions shared across images lower the loss; the construction is
the sample splitting of double machine learning \citep{chernozhukov2018dml} applied to a training loss.
The held-out explained variance of the batch fit is logged as a memorization-free diagnostic.

\paragraph{First version: failure.} The first runs produced NaN features on three of four Spur-CIFAR10
runs and negative held-out explained variance everywhere. The ridge scaled with the feature norm and was
small, so fits of 64 images in 512 dimensions nearly interpolated and the solve became ill conditioned.
Scaling every feature row to unit norm and using an absolute ridge $\rho=1$ keeps the eigenvalues of
the system at least $\rho$ \citep{hoerl1970ridge}.

\begin{table}[h]
\centering
\caption{Repaired cross-fit, seeds 1 and 2. Batch EV is the held-out explained variance of the batch fit
at epoch 500.}
\label{tab:xfit}
\small
\begin{tabular}{l l r r r}
\toprule
Dataset & Arm & WGA & Acc. & Batch EV \\
\midrule
Spur-CIFAR10 & SimCLR & $15.6\pm1.0$ & 59.5 & \\
 & F2, standardized & $32.9\pm1.6$ & 69.5 & 0.10 \\
 & xfit & $33.6\pm2.3$ & 69.7 & 0.13 \\
 & xfit, shuffled & $18.2\pm3.2$ & 57.5 & $-0.02$ \\
\midrule
MetaShift & xfit & $43.9\pm3.2$ & 54.8 & $-0.01$ \\
 & xfit, shuffled & $41.7\pm0.0$ & 54.4 & $-0.03$ \\
 & xfit, ridge 10 & $42.5\pm2.2$ & 54.4 & $-0.01$ \\
 & xfit, weight 3 & $43.9\pm0.7$ & 57.8 & $0.00$ \\
 & xfit, balancing weights & $46.9\pm4.4$ & 58.5 & $-0.01$ \\
\bottomrule
\end{tabular}
\end{table}

On Spur-CIFAR10 the cross-fit matches F2, and its shuffled control falls to the SimCLR level
(Table~\ref{tab:xfit}). Where the shuffled F2 had kept nine points, the cross-fitted control keeps
2.6, so the whole remaining gain comes from the concept content. We counted this as the first clean
attribution of the branch. On MetaShift the 145 factors stayed at the noise level under every setting,
which pointed away from the loss and toward the factors themselves.

\section{Concept grouping revisited}
\label{sec:grouping}

Inspecting single images showed that the factors were rarely groups at all. On Waterbirds with Open
Images at the default thresholds, 499 of 502 groups were single concepts, and the strongest concepts of
an image often fell below the frequency cut (``Footpath'' on 0.94\% of the images). We measured four
candidate similarity signals on concept pairs of the Waterbirds LAION cache (Table~\ref{tab:signals}).

\begin{table}[h]
\centering
\caption{Similarity of concept pairs on Waterbirds (LAION). Text: cosine of the centred SpLiCE
dictionary directions; raw text: cosine of the uncentred CLIP text embeddings; co-activation: cosine
of the code columns; response: correlation of the images' CLIP alignment with the two concepts.}
\label{tab:signals}
\small
\begin{tabular}{l r r r r}
\toprule
Pair & Text & Raw text & Co-activation & Response \\
\midrule
path, walkway & 0.23 & 0.85 & 0.37 & 0.66 \\
heron, egret & 0.70 & 0.87 & 0.09 & 0.79 \\
ocean, sea & 0.71 & 0.94 & 0.00 & 0.94 \\
forest, woodland & 0.49 & 0.90 & 0.08 & 0.89 \\
\midrule
bamboo, jungle & 0.24 & 0.67 & 0.10 & 0.67 \\
gull, seascape & 0.16 & 0.65 & 0.02 & 0.58 \\
heron, marsh & 0.08 & 0.67 & 0.06 & 0.11 \\
\bottomrule
\end{tabular}
\end{table}

Two properties of SpLiCE explain the failure of the original rule, which links two concepts only when
both their centred text similarity and their co-activation pass a threshold. The L1 penalty selects one
of several correlated dictionary atoms, the known behaviour of the lasso on correlated predictors
\citep{tibshirani1996lasso,zou2005elasticnet}, so synonyms rarely fire on the same image. Centring the
dictionary removes the shared component of related words and lowers their cosine. Lowering both
thresholds then joins concepts through chains of pairwise links, the single-linkage behaviour of
connected components: at text 0.3 and co-activation 0.1 one group held hummingbird, parrot, jungle and
bamboo, so the factor itself fused a bird with its background.

The \emph{meaning grouping} clusters concepts by average linkage \citep{sokal1958upgma} on the raw CLIP
text embeddings, where synonyms reach 0.85 to 0.94 and object--context pairs 0.64 to 0.72. Average
linkage joins two clusters only when their members agree on average, which stops the chains. An
optional veto drops pairs whose image responses correlate below 0.5. Concepts from 0.2\% frequency
enter the grouping, so rare synonyms join before the factor band applies to their sum. A sweep over
576 factor sets on MetaShift and Waterbirds chose text 0.85 and response 0.5 on MetaShift, which yields
groups such as \{kitty, kitten, cat\}, \{canine, dog\} and \{cozy, comfy, cosy\} with LAION. Open
Images groups also join homonyms and similar Latin names (Quail and Quail meat).

The same sweep put pair precision, the share of the top correlated pairs that join a class factor to an
attribute factor, at zero for both grouping methods. The phi coefficient ranks within-scene pairs first
(windshield and driving 0.64, toilet and sink 0.57) and class--context pairs last (dog and lawn 0.18,
cat and couch 0.04), because phi penalizes pairs of very different frequency (cat on 36\% of the images,
couch on 3\%). We set F1 aside: F2 and xfit use every factor and need no pairs.

\section{Learnability: what the student can and does learn}
\label{sec:learnability}

\subsection{A bound from CLIP features}
The cross-fit can only reward what a ridge fitted on 64 images can predict. We therefore ran the same
fit on frozen CLIP image embeddings, the features the targets come from, as an upper bound
(Table~\ref{tab:bound}). With the 278 meaning factors from 1\%, even CLIP scores below zero: a factor on
1\% of the images occurs 0.6 times per half batch, and the loss consists of noise. Restricting the
factors to those on at least 5\% of the images (24 factors) raises the bound to 0.16 against $-0.07$
for shuffled targets.

\begin{table}[h]
\centering
\caption{Held-out explained variance of a ridge ($\rho=1$) fitted on $n$ MetaShift images and evaluated
on $n$ others, averaged over 40 draws.}
\label{tab:bound}
\small
\begin{tabular}{l r r r}
\toprule
Factors & $n$ & CLIP features & Shuffled targets \\
\midrule
278 (from 1\%) & 64 & $-0.015$ & $-0.079$ \\
278 (from 1\%) & 800 & $0.140$ & $-0.086$ \\
24 (from 5\%) & 64 & $0.161$ & $-0.069$ \\
11 (from 10\%) & 64 & $0.218$ & $-0.061$ \\
\bottomrule
\end{tabular}
\end{table}

With 24 factors and weight 1 the student's batch score rose to 0.05 at epoch 500 and was still rising,
against $-0.02$ for the control, while WGA stayed at $46.0\pm2.2$ against $46.7\pm1.7$. We read the
rising score as concepts being learned and raised the weight.

\subsection{Weight 3 and 10: scores above the bound}
Every concept-factor run now writes a per-factor record: a five-fold ridge ($\rho=10$) from the
training features to each real standardized factor, beside the score of the same fit on CLIP features.
Shuffled arms score the real factors too, so they show what SimCLR learns of the concepts by itself.

\begin{table}[h]
\centering
\caption{MetaShift, 24 meaning factors, seeds 1 and 2. Train EV is the five-fold score on the training
images; CLIP features score 0.26.}
\label{tab:weights}
\small
\begin{tabular}{l r r r c}
\toprule
Arm & Train EV & WGA & Acc. & Groups \\
\midrule
xfit, weight 3, shuffled & 0.02 & 47.6 & 58.1 & 68 / 53 / 50 / 59 \\
xfit, weight 3 & 0.33 & 43.2 & 59.6 & 74 / 45 / 48 / 69 \\
xfit, weight 10 & 0.79 & 45.5 & 61.0 & 74 / 51 / 48 / 71 \\
\bottomrule
\end{tabular}
\end{table}

At weight 10 every factor reached 0.72 to 0.87 on the training images, three times the CLIP score and
alike for factors CLIP explains well (cat, 0.68) and badly (helper, 0.14). A logistic probe on the 24
factor values themselves reaches 77\% WGA and 89\% accuracy on group-balanced splits of the training
images, and CLIP embeddings reach 87\% and 92\%; the student stays at 45\% and 61\% on validation. The
factors carry the information the probe needs; the student holds something else.

\subsection{Unseen images: memorization in the backbone}
Cross-fitting holds images out of the ridge and leaves the backbone free to place each training image's
target in its own features, where one shared linear map reads it back; the whole-dataset score cannot
see this. We kept a fixed fifth of the training images out of the concept loss (SimCLR still trains on
them) and scored the factors on those images with a ridge fitted on the others.

\begin{table}[h]
\centering
\caption{Explained variance of the factors on the images in the concept loss (seen) and on the fifth
held out from it (unseen), epoch 500, seeds 1 and 2.}
\label{tab:unseen}
\small
\begin{tabular}{l l r r r r r}
\toprule
Dataset & Arm & Seen & Unseen & CLIP unseen & WGA & Acc. \\
\midrule
MetaShift & xfit, weight 10 & 0.60 & 0.04 & 0.26 & 47.9 & 60.4 \\
 & shuffled & 0.01 & 0.02 & 0.26 & 43.8 & 52.6 \\
Spur-CIFAR10 & xfit & 0.31 & 0.27 & 0.42 & 30.2 & 68.7 \\
 & shuffled & 0.10 & 0.10 & 0.42 & 21.7 & 58.1 \\
\bottomrule
\end{tabular}
\end{table}

Table~\ref{tab:unseen} separates the two datasets. On Spur-CIFAR10 the student carries the factors to new
images at two thirds of the CLIP score, for the classes (aircraft 0.75, horse 0.74, truck 0.72) and for
the line colour (purple 0.62) alike, the separate class and attribute directions that deep feature
reweighting needs. On MetaShift the unseen score peaks near 0.08 at epochs 100 to 200 and falls to 0.04
while the seen score climbs to 0.60; the class factors carry over in part (cat 0.37, dog 0.23) and the
context factors do not (cozy 0.03). The student memorizes the targets of MetaShift's 1,700 images and
generalizes them on the 45,000 images of Spur-CIFAR10. This explains why the majority groups of
MetaShift rose in every round while the minority groups stayed flat.

\section{Waterbirds}
\label{sec:waterbirds}

Waterbirds received frozen configurations throughout. F2 as frozen after round 1 lowered WGA to
$41.6\pm6.1$ against $45.8\pm1.7$ for SimCLR over four seeds. The repaired cross-fit with the default
217 Open Images factors learned little and memorized nothing (Table~\ref{tab:waterbirds}): seen and
unseen scores agree and exceed the control by a few hundredths, and validation accuracy stays at chance.

\begin{table}[h]
\centering
\caption{Waterbirds validation. SimCLR: seeds 1 to 4; other arms: seeds 1 and 2 with a fifth of the
images held out of the concept loss; the meaning arms report the final probe.}
\label{tab:waterbirds}
\small
\begin{tabular}{l r r r r c}
\toprule
Arm & Seen & Unseen & WGA & Acc. & Groups \\
\midrule
SimCLR & & & $45.8\pm1.7$ & 52.3 & 47 / 54 / 52 / 64 \\
xfit, 217 factors, weight 1 & 0.04 & 0.03 & 46.5 & 51.1 & \\
\quad shuffled & 0.02 & 0.02 & 43.9 & 50.5 & \\
xfit, 70 meaning factors, weight 10 & 0.30 & 0.09 & \textbf{55.3} & \textbf{59.2} & 56 / 59 / 66 / 68 \\
\quad shuffled & 0.01 & 0.02 & 42.8 & 52.6 & 43 / 57 / 58 / 66 \\
\bottomrule
\end{tabular}
\end{table}

The configuration that came out of the MetaShift diagnostics, meaning groups at text 0.85 and
response 0.5, factors from 5\% and weight 10, lifts every Waterbirds group, minority groups included,
to 55.3\% WGA (57.8 and 52.7 for the two seeds). This exceeds SimCLR by 9.5 points, its shuffled control
by 12.5 points and the CoSpRo validation result ($50.0\pm3.2$) by 5.3 points. The student stores part of
the targets (seen 0.30) and carries a third of the CLIP score to unseen images (0.09 against 0.02 for
the control), more than twice the MetaShift value. Waterbirds, with 4,795 images, sits between the
memorizing and the generalizing dataset.

We count this as the first success of the branch on a dataset where CoSpRo already worked, with two
caveats: two seeds, and the hold-out of a fifth of the images. Its settings came from MetaShift
diagnostics and the learnability bounds; the weight of 10 was chosen to probe memorization on
MetaShift, and no setting was tuned on Waterbirds labels. The configuration is frozen from here on.

\section{Where we stand and how we continue}
\label{sec:next}

\paragraph{Conclusions so far.}
\begin{enumerate}
\item Distilling SpLiCE concept factors into a SimCLR student improves worst-group accuracy when the
student can generalize the factors: Spur-CIFAR10 (WGA 15.6 to 33.6) and Waterbirds (45.8 to 55.3).
\item Cross-fitting makes the attribution clean: shuffled targets lose the gain, so the concept content
carries it.
\item On small datasets an auxiliary per-image target is memorized. Cross-fitting removes memorization in
the head; the backbone still memorizes, and only a score on images outside the concept loss shows it.
\item The grouping must join synonyms that SpLiCE never activates together and must avoid chains through
context; the meaning grouping does both. Factors must be frequent enough for a half batch to contain
them.
\item Pair selection by phi favours within-scene pairs, so conditioned batches (F1) need another pair
statistic before they can work.
\end{enumerate}

\paragraph{Running and planned experiments.}
CelebA trains SimCLR, the weight-1 cross-fit and its control for 250 epochs on 162,770 images; at epoch
188 the cross-fit scores 0.20 on the training factors against 0.07 for the control and 0.28 for CLIP.
Waterbirds runs the frozen configuration on seeds 3 and 4 with the hold-out and on seeds 1 to 4 without
it, which gives the headline number. Spur-CIFAR10 runs the same configuration to test that one setting
serves all datasets, and CelebA follows with it once the weight-1 runs finish. At most eight jobs run at
a time.

\paragraph{Open questions.}
Why does Waterbirds generalize part of the factors where MetaShift does not, beyond its threefold size?
The 70 Open Images factors include bird families and backgrounds with high CLIP scores (Water bird 0.61,
Bamboo 0.76), while the MetaShift factors are dominated by pet-photo vocabulary. Can memorization be
made unprofitable on small datasets, for example by mixing images and their targets as in mixup
\citep{zhang2018mixup}, so that storing per-image targets stops lowering the loss? Does a pair statistic
that measures how much a context raises the probability of a class revive F1?

\paragraph{Threats to validity.}
Two seeds per arm leave wide intervals, and MetaShift's 72-image validation groups make its WGA noisy.
All values are validation results; the locked test split stays unused until the final configuration is
fixed. Waterbirds served as a diagnostic dataset for the weight-1 cross-fit before the configuration was
frozen, so its role as a held-out dataset rests on the absence of any label-based selection on it.

\appendix
\section{Options and arms}
\label{app:options}

\begin{table}[h]
\centering
\caption{Training options of the concept-factor method (\texttt{--splice\_mode concept\_factors}).}
\small
\begin{tabularx}{\linewidth}{l X}
\toprule
Option & Meaning \\
\midrule
\texttt{--factor\_distill\_weight} & F2 or xfit loss weight $\lambda$ \\
\texttt{--factor\_targets} & standardized, whitened, presence or shuffled \\
\texttt{--factor\_cross\_fit}, \texttt{--factor\_ridge} & cross-fitted ridge and its ridge $\rho$ \\
\texttt{--factor\_sample\_weighting} & none, atypicality or balanced \\
\texttt{--factor\_min\_frequency} & lowest factor frequency \\
\texttt{--factor\_merge\_similarity} & merge groups by image response; 0 disables \\
\texttt{--factor\_concept\_groups} & groups file, e.g.\ meaning groups from \texttt{build\_meaning\_groups} \\
\texttt{--factor\_holdout\_fraction} & share of images kept out of the concept loss \\
\texttt{--factor\_condition\_fraction} & F1 share of conditioned batches \\
\texttt{--factor\_block\_weight} & F3 concept-block loss weight \\
\bottomrule
\end{tabularx}
\end{table}

Run records live under \texttt{outputs/seeds/factors\_<dataset>/}, study summaries under
\texttt{outputs/reports/factors\_<dataset>/summary.md} and per-factor learnability records in
\texttt{factor\_learnability\_epoch\_<E>.json} beside each run's status file. The arms are defined in
\texttt{scripts/submit\_factor\_study.sh}; the running notes are in
\texttt{docs/notes/CONCEPT\_FACTORS\_2026\_09\_26.md}.

\begingroup
\footnotesize
\setlength{\bibsep}{3pt plus 1pt}
\begin{thebibliography}{99}

\bibitem[Arpit et~al.(2017)]{arpit2017closer}
D.~Arpit, S.~Jastrz\k{e}bski, N.~Ballas, D.~Krueger, E.~Bengio, M.~S.~Kanwal, T.~Maharaj, A.~Fischer,
A.~Courville, Y.~Bengio and S.~Lacoste-Julien.
A Closer Look at Memorization in Deep Networks. \emph{ICML}, 2017.
\url{https://arxiv.org/abs/1706.05394}.

\bibitem[Bertinetto et~al.(2019)]{bertinetto2019r2d2}
L.~Bertinetto, J.~F.~Henriques, P.~H.~S.~Torr and A.~Vedaldi.
Meta-learning with differentiable closed-form solvers. \emph{ICLR}, 2019.
\url{https://arxiv.org/abs/1805.08136}.

\bibitem[Bhalla et~al.(2024)]{bhalla2024splice}
U.~Bhalla, A.~Oesterling, S.~Srinivas, F.~P.~Calmon and H.~Lakkaraju.
Interpreting CLIP with Sparse Linear Concept Embeddings (SpLiCE). \emph{NeurIPS}, 2024.
\url{https://arxiv.org/abs/2402.10376}.

\bibitem[Chen et~al.(2020)]{chen2020simclr}
T.~Chen, S.~Kornblith, M.~Norouzi and G.~Hinton.
A Simple Framework for Contrastive Learning of Visual Representations. \emph{ICML}, 2020.
\url{https://arxiv.org/abs/2002.05709}.

\bibitem[Chen et~al.(2021)]{chen2021intriguing}
T.~Chen, C.~Luo and L.~Li.
Intriguing Properties of Contrastive Losses. \emph{NeurIPS}, 2021.
\url{https://arxiv.org/abs/2011.02803}.

\bibitem[Chernozhukov et~al.(2018)]{chernozhukov2018dml}
V.~Chernozhukov, D.~Chetverikov, M.~Demirer, E.~Duflo, C.~Hansen, W.~Newey and J.~Robins.
Double/debiased machine learning for treatment and structural parameters.
\emph{The Econometrics Journal}, 21(1):C1--C68, 2018.

\bibitem[Hamidieh et~al.(2024)]{hamidieh2024views}
K.~Hamidieh, H.~Zhang, S.~Sankaranarayanan and M.~Ghassemi.
Views Can Be Deceiving: Improved SSL Through Feature Space Augmentation. \emph{ICLR}, 2024.

\bibitem[Hoerl and Kennard(1970)]{hoerl1970ridge}
A.~E.~Hoerl and R.~W.~Kennard.
Ridge Regression: Biased Estimation for Nonorthogonal Problems. \emph{Technometrics}, 12(1):55--67, 1970.

\bibitem[Kessy et~al.(2018)]{kessy2018whitening}
A.~Kessy, A.~Lewin and K.~Strimmer.
Optimal Whitening and Decorrelation. \emph{The American Statistician}, 72(4):309--314, 2018.

\bibitem[Khosla et~al.(2020)]{khosla2020supcon}
P.~Khosla et~al.
Supervised Contrastive Learning. \emph{NeurIPS}, 2020.
\url{https://arxiv.org/abs/2004.11362}.

\bibitem[Kirichenko et~al.(2023)]{kirichenko2023dfr}
P.~Kirichenko, P.~Izmailov and A.~G.~Wilson.
Last Layer Re-Training is Sufficient for Robustness to Spurious Correlations. \emph{ICLR}, 2023.
\url{https://arxiv.org/abs/2204.02937}.

\bibitem[Lee et~al.(2019)]{lee2019metaoptnet}
K.~Lee, S.~Maji, A.~Ravichandran and S.~Soatto.
Meta-Learning with Differentiable Convex Optimization. \emph{CVPR}, 2019.
\url{https://arxiv.org/abs/1904.03758}.

\bibitem[Liang and Zou(2022)]{liang2022metashift}
W.~Liang and J.~Zou.
MetaShift: A Dataset of Datasets for Evaluating Contextual Distribution Shifts and Training Conflicts.
\emph{ICLR}, 2022. \url{https://arxiv.org/abs/2202.06523}.

\bibitem[Liu et~al.(2015)]{liu2015celeba}
Z.~Liu, P.~Luo, X.~Wang and X.~Tang.
Deep Learning Face Attributes in the Wild. \emph{ICCV}, 2015.

\bibitem[Locatello et~al.(2020)]{locatello2020weakly}
F.~Locatello, B.~Poole, G.~R\"atsch, B.~Sch\"olkopf, O.~Bachem and M.~Tschannen.
Weakly-Supervised Disentanglement Without Compromises. \emph{ICML}, 2020.
\url{https://arxiv.org/abs/2002.02886}.

\bibitem[Radford et~al.(2021)]{radford2021clip}
A.~Radford et~al.
Learning Transferable Visual Models From Natural Language Supervision. \emph{ICML}, 2021.
\url{https://arxiv.org/abs/2103.00020}.

\bibitem[Robinson et~al.(2021)]{robinson2021shortcut}
J.~Robinson, L.~Sun, K.~Yu, K.~Batmanghelich, S.~Jegelka and S.~Sra.
Can contrastive learning avoid shortcut solutions? \emph{NeurIPS}, 2021.
\url{https://arxiv.org/abs/2106.11230}.

\bibitem[Sagawa et~al.(2020)]{sagawa2020groupdro}
S.~Sagawa, P.~W.~Koh, T.~B.~Hashimoto and P.~Liang.
Distributionally Robust Neural Networks for Group Shifts: On the Importance of Regularization for
Worst-Case Generalization. \emph{ICLR}, 2020. \url{https://arxiv.org/abs/1911.08731}.

\bibitem[Sokal and Michener(1958)]{sokal1958upgma}
R.~R.~Sokal and C.~D.~Michener.
A Statistical Method for Evaluating Systematic Relationships.
\emph{University of Kansas Science Bulletin}, 38:1409--1438, 1958.

\bibitem[Tibshirani(1996)]{tibshirani1996lasso}
R.~Tibshirani.
Regression Shrinkage and Selection via the Lasso.
\emph{Journal of the Royal Statistical Society B}, 58(1):267--288, 1996.

\bibitem[Tsai et~al.(2021)]{tsai2021conditional}
Y.-H.~H.~Tsai, M.~Q.~Ma, H.~Zhao, K.~Zhang, L.-P.~Morency and R.~Salakhutdinov.
Conditional Contrastive Learning for Improving Fairness in Self-Supervised Learning. 2021.
\url{https://arxiv.org/abs/2106.02866}.

\bibitem[Yao et~al.(2024)]{yao2024multiview}
D.~Yao, D.~Xu, S.~Lachapelle, S.~Magliacane, P.~Taslakian, G.~Martius, J.~von K\"ugelgen and F.~Locatello.
Multi-View Causal Representation Learning with Partial Observability. \emph{ICLR}, 2024.
\url{https://arxiv.org/abs/2311.04056}.

\bibitem[Zhang et~al.(2017)]{zhang2017rethinking}
C.~Zhang, S.~Bengio, M.~Hardt, B.~Recht and O.~Vinyals.
Understanding Deep Learning Requires Rethinking Generalization. \emph{ICLR}, 2017.
\url{https://arxiv.org/abs/1611.03530}.

\bibitem[Zhang et~al.(2018)]{zhang2018mixup}
H.~Zhang, M.~Cisse, Y.~N.~Dauphin and D.~Lopez-Paz.
mixup: Beyond Empirical Risk Minimization. \emph{ICLR}, 2018.
\url{https://arxiv.org/abs/1710.09412}.

\bibitem[Zhang et~al.(2021)]{zhang2021stable}
X.~Zhang, P.~Cui, R.~Xu, L.~Zhou, Y.~He and Z.~Shen.
Deep Stable Learning for Out-of-Distribution Generalization. \emph{CVPR}, 2021.
\url{https://arxiv.org/abs/2104.07876}.

\bibitem[Zhu et~al.(2023)]{zhu2023lassl}
W.~Zhu, S.~Liu, C.~Fernandez-Granda and N.~Razavian.
Making Self-supervised Learning Robust to Spurious Correlation via Learning-speed Aware Sampling.
\emph{arXiv:2311.16361}, 2023. \url{https://arxiv.org/abs/2311.16361}.

\bibitem[Zou and Hastie(2005)]{zou2005elasticnet}
H.~Zou and T.~Hastie.
Regularization and Variable Selection via the Elastic Net.
\emph{Journal of the Royal Statistical Society B}, 67(2):301--320, 2005.

\end{thebibliography}
\endgroup

\end{document}
